% called by main.tex
%

\section*{Abstract}
\addcontentsline{toc}{section}{Abstract}
\label{sec::abstract}

One of the major problems in modern society is road traffic accidents, which represent one of the leading causes of death among young people aged 5 to 29. Among their causes, driver distraction and drowsiness are major contributors and can occur for various reasons, some of which are beyond human control; however, signs of distraction or drowsiness can be detected through different indicators and physiological signals. Currently, several systems can detect these conditions in vehicle drivers to determine whether they are experiencing drowsiness or distraction, but they tend to be quite robust and may require expensive or sometimes hard-to-obtain hardware.

Given this problem, our general objective is to develop an integrated camera system capable of detecting distraction and/or drowsiness in a driver in order to alert both the user and a third party, while remaining lightweight enough to be integrated into a single-board computer that is more affordable and accessible than current options. Alongside this, our specific objectives are to develop classifiers to detect drowsiness and distraction. The system was integrated into a single-board computer with a webcam and audio alerts to operate in real time. Finally, a system was integrated to send notifications to a third party through a Telegram bot.

To detect drowsiness, our system implemented blink duration and PERCLOS detection using the 68-point facial landmark detector from Dlib to extract facial features and compute the Eye Aspect Ratio, which reflects eye behavior, in order to measure how long the eyelid remains closed and compute PERCLOS over a 30-second window. Another way drowsiness episodes were detected was through yawn detection, using a custom convolutional neural network trained on the YawDD dataset. The last way drowsiness episodes were detected was through head pitch angle, calculated from the facial landmarks detected by Dlib to estimate the angle of the user's head.

To detect user distraction, several neural network architectures were trained using the DMD dataset, leading to the decision to use a C-LSTM in the final implementation. This C-LSTM showed the best results among all experimented neural network architectures in terms of distraction event detection capability, although it was also considerably heavier than the other architectures; however, through certain optimizations, its execution on the selected single-board computer was achieved.

Finally, for implementing all developed systems on a single-board computer, a Raspberry Pi 5B was chosen, and to ensure its execution in the target environment, several code optimizations were carried out: first, the code was modularized to make modifications easier, and finally, the program was adapted to C++ as a form of optimization.

During tests with the datasets, the program showed promising results across all areas, showing only minor deficiencies in false positive detections within blink detection. During the project, certain improvements were identified for future work, such as the power supply system for the device inside the vehicle and reducing false positives in blink detection when tilting the head downward. Despite this, real-time tests demonstrated the expected execution of the code as a whole, including all its somnolence and distraction detection features, and others such as audio notifications and messages via Telegram. 


\newpage
\section*{Resumen}
\addcontentsline{toc}{section}{Resumen}
\label{sec::resumen}


Una de las mayores problemáticas presentes en la sociedad moderna es la de los accidentes viales, los cuales representan una de las mayores causas de muerte en jóvenes de entre los 5 a 29 años. Hablando de las causas de estos, la distracción y la somnolencia al volante se encuentran entre las principales causas de accidentes viales, y estas dos pueden suceder por diversas razones, las cuales pueden estar fuera del control de las personas, pero fuera de esto, existe la capacidad de poder detectar señales de distracción o somnolencia a través de diversas señales fisiológicas. Actualmente, existen diversos sistemas capaces de detectar señales fisiológicas en los conductores de vehículos para determinar si sufren algún episodio de somnolencia o distracción, pero estos suelen ser sistemas bastante robustos, y pueden llegar a requerir hardware costoso y en ocasiones difícil de obtener.

Conociendo esta problemática, tenemos propuesto como objetivo general el desarrollar un sistema integrado con cámara que sea capaz de detectar distracción y/o somnolencia en un conductor para poder alertar al usuario y a un tercero, y que este sea más ligero como para poder integrarlo en una SBC que pueda ser más económica y accesible que las opciones actuales. Junto con este, tenemos los objetivos específicos de desarrollar clasificadores para detectar somnolencia y distracción. Se integró el sistema en una SBC con cámara web y alertas de audio para operar en tiempo real. Finalmente, se integró un sistema capaz de enviar notificaciones a un tercero a través de un bot de Telegram.

Para poder detectar somnolencia, se implementó en nuestro sistema la funcionalidad de detectar duración de parpadeos y PERCLOS utilizando el detector de 68 puntos de Dlib, capaz de extraer características del rostro para poder utilizarlas y obtener la medida del EAR, que muestra el comportamiento del ojo, para poder así detectar cuanto tiempo permanece cerrado el párpado y el PERCLOS dentro de una ventana de 30 segundos. Otra de las formas en las que se detectaron episodios de somnolencia fue la detección de bostezos, estos siendo detectados a través de una CNN propia entrenada con el conjunto de datos de YawDD. La última forma con la que se detectaron episodios de somnolencia fue el ángulo de cabeceo, este se calculó a través de los puntos detectados por Dlib para estimar el ángulo en el que se encuentra la cabeza del usuario.

Para poder detectar la distracción del usuario, se entrenaron diversas estructuras de redes neuronales utilizando el conjunto de datos de DMD, terminando por la decisión de utilizar una C-LSTM en la implementación final. Esta C-LSTM mostró los mejores resultados de todas las arquitecturas de redes neuronales en las que se experimentó al hablar de su capacidad de detección de eventos de distracción, pero a su vez fue considerablemente más pesada que el resto de las arquitecturas, aunque a través de ciertas optimizaciones, se logró su ejecución en la SBC elegida.

Finalmente, para la implementación de todos los sistemas desarrollados en una SBC, se eligió utilizar una Raspberry Pi 5B, y para poder asegurar su ejecución dentro del entorno deseado, se llegaron a realizar diversas optimizaciones al código, primero, se realizó una modularización del código para poder realizar cambios de forma más sencilla, y finalmente se realizó una adaptación del programa a C++ como forma de optimización.

Durante pruebas con los conjuntos de datos, el programa mostró resultados prometedores en todas las áreas, únicamente mostrando leves deficiencias al momento de la detección de falsos positivos dentro de la detección de parpadeos. Dentro del proyecto se llegaron a identificar ciertas mejoras para trabajo a futuro, como el sistema de alimentación del equipo dentro del vehículo y la reducción de falsos positivos en la detección de parpadeos al inclinar la cabeza hacia abajo. A pesar de lo anterior, las pruebas en tiempo real demostraron la ejecución esperada del código en su totalidad, incluyendo todas sus funciones de detección de somnolencia y distracción, así como el sistema de notificaciones por audio y por mensaje vía Telegram. 
